\documentclass{article}
\usepackage{graphicx} 
\usepackage{xcolor}
\usepackage{amsmath, amssymb}
% Bibliography 
\usepackage[backend=biber,style=authoryear]{biblatex}
\addbibresource{references.bib} 

\begin{document}

%=================================================================
% 1) Deckblatt
%-----------------------------------------------------------------
\begin{titlepage}
  \vspace*{2cm}
  \begin{center}
    % Logo
    \includegraphics[width=0.7\textwidth]{Bildmarke_black_23cm.png}\\[2cm]

    {\LARGE\bfseries Projekt 10: Bodenatmung\par}

    \vspace{2cm}

    \large
    Friedrich-Schiller-Universität Jena\\
    Fakultät für Mathematik und Informatik\\
    Statistische Verfahren WW 2025

    \vfill

    \Large
    Anshul Khubchandani et al. 

    \vspace{1.5cm}

    \normalsize
  \end{center}
\end{titlepage}

% ---------------------------------------------------------
\tableofcontents
\thispagestyle{empty}
\addtocounter{page}{-1}

% ---------------------------------------------------------
\newpage
\section{Organisatorisches}
\textcolor{red}{Abgabetermin: 31.03.2026}


Umfang:


• keine explizite Vorgabe

• Aufgabenstellung muss vollständig bearbeitet sein

• 15 Seiten reichen sicher aus \newline

Abgabe:

• Projektarbeit in elektronischer Form

• Selbstständigkeitserklärung (keine Zulassung von KI!)

• zulässige Sprachen: deutsch, englisch (ungarisch)

• R-Skripte (werden nicht bewertet, aber zur Kontrolle benötigt)\newline 

Diese Seite wird am ende gelöscht :)

% ---------------------------------------------------------
\newpage
\section{Einleitung}
\textcolor{red} {• von der Fragestellung zur Statistik
• kurze Beschreibung der Fragestellung
• Hinführen zur statistischen Modellierung
• warum statistische Modellierung
• welche Modellklasse}

Die Bodenatmung beschreibt die Freisetzung von Kohlenstoffdioxid (CO$_2$) aus dem Boden und bildet den 
zweitgrößten Kohlenstofffluss in terrestrischen Ökosystemen nach der Photosynthese. \cite{raich1992global} Sie setzt
sich aus der mikrobielle Aktivität und der Wurzelatmung zusammen und bildet dafür zentrale Maße. Der Anteil 
der Wurzelatmung schwankt je nach Ökosystem und Jahreszeit zwischen 10\% und 90\%. \cite{Hanson2000}
Die Bodenatmung stellt einen wichtigen Bestandteil des globalen Kohlenstoffkreislaufs dar. Böden sind das
größte terrestrische Reservoir für organischen Kohlenstoff und speichern etwa 71\% bis 81\% des gesamten
terrestrischen Kohlenstoffs. \cite{IPCC2001} Die Bodenatmung ist der Hauptweg auf dem Kohlenstoff wieder in die 
Atmosphäre zurückkehrt, wodurch kleine Änderungen in dessen Intensität große Auswirkungen auf die globale 
CO$_2$-Konzentration und damit auf das Klima hat. \cite{1997biogeochemistry,IPCC2001,Churkina2003}

Im Nationalpark Hainich wurden umfangreiche Messungen zur Bodenatmung sowie zu physikalischen, chemischen und
biologischen Bodenparametern durchgeführt. Diese Daten ermöglichen die statistische Modellierung der Bodenatmung
in Abhängigkeit von Umweltvariablen wie Bodentemperatur, Bodenfeuchte, pH-Wert und Nährstoffgehalten. Das Ziel
der Forschung ist es, die Sensitivität der Bodenatmung gegenüber dem Klimawandel besser vorherzusagen. \cite{diss}

Ziel dieser Arbeit ist die Herleitung eines einfachen, aber aussagekräftigen Prognosemodells für die
Bodenatmung auf Basis des gegebenen Datensatzes. Dies wird mit einem klassisch linearem Modell auf
log-transformierten Daten umgesetzt. Darüber hinaus wird im Rahmen einer Simulationsstudie
untersucht, welche Auswirkungen eine falsche Anwendung von t-Tests im Kontext der Vorwärtsselektion von
Prädiktoren auf die Irrtumswahrscheinlichkeit hat.
% ---------------------------------------------------------
\newpage
\section{Material und Methoden}
\textcolor{red}{• Beschreibung des Datensatzes • statistisches Modell!!!
• deterministischer und stochastischer Teil
• verwendete statistische Methoden
• Beschreibung der Vorgehensweise
• kein Beweis theoretischer Resultate, aber Definition der verwendeten Größen
• nachvollziehbare Beschreibung der Simulation}

\subsection{Datensatz}

Die Analyse basiert auf dem Datensatz \texttt{hainich.csv}, der von A.~Soe (MPI Biogeochemie)
bereitgestellt wurde. Der Datensatz enthält 38 Messungen der Bodenatmung (\texttt{soil.res}) sowie
folgende Einflussgrößen:

\begin{itemize}
\setlength{\itemsep}{0pt}
\item \texttt{temp.0}, \texttt{temp.5}, \texttt{temp.10}, \texttt{temp.15} – Bodentemperatur in 0, 5, 10, 15 cm Tiefe
\item \texttt{smoi} – relative Bodenfeuchte (soil moisture)
\item \texttt{lmoi} – relative Feuchte der Streuschicht (litter moisture)
\item \texttt{litter.d} – Dicke der Streuschicht
\item \texttt{ph.litter}, \texttt{ph.soil0}, \texttt{ph.soil5} – pH-Wert der Streuschicht bzw. des Bodens in 0-5 cm bzw. 5-10 cm Tiefe
\item \texttt{rootdw0}, \texttt{rootdw5} – Trockenmasse der Wurzeln in 0-5 cm bzw. 5-10 cm Tiefe
\item \texttt{rootn0}, \texttt{rootn5} – Stickstoffgehalt der Wurzeln in 0-5 cm bzw. 5-10 cm Tiefe
\item \texttt{rootc0}, \texttt{rootc5} – Kohlenstoffgehalt der Wurzeln in 0-5 cm bzw. 5-10 cm Tiefe
\item \texttt{litter.c}, \texttt{litter.n} – Kohlenstoff- bzw. Stickstoffgehalt der Streuschicht
\item \texttt{soiln0}, \texttt{soiln5} – Stickstoffgehalt des Bodens in 0-5 cm bzw. 5-10 cm Tiefe
\item \texttt{soilc0}, \texttt{soilc5} – Kohlenstoffgehalt der Bodens in 0-5 cm bzw. 5-10 cm Tiefe
\item \texttt{litno3}, \texttt{sno30}, \texttt{sno35} – Nitratgehalt der Streuschicht bzw. des Bodens in 0-5 cm bzw. 5-10 cm Tiefe
\item \texttt{litnh4}, \texttt{snh40}, \texttt{snh45} - Ammoniumgehalt der Streuschicht bzw. des Bodens in 0-5 cm bzw. 5-10 cm Tiefe
\item \texttt{litdoc}, \texttt{sdoc0}, \texttt{sdoc5} – Gehalt an gelöstem organischen Kohlenstoff (DOC - dissolved organic carbon) der Streuschicht bzw. des Bodens in 0-5 cm bzw. 5-10 cm Tiefe
\item \texttt{nn} – Entfernung der Probestelle zum nächsten Baum (nearest neighbour)
\end{itemize}

\subsubsection{Transformation}
\subsubsection{Korrelation}
\subsubsection{Interaktionen zwischen den Einflussgrößen}

folgende Interaktionen gelten als biologisch relevant:
\texttt{lmoi} und \texttt{litter.de}:
Je dicker die Streuschicht ist, desto mehr Wasser kann sie speichern und desto länger dauert es,
bis sie ausgetrocknet ist.

\texttt{temp.15} und \texttt{lmoi}:
Mikroorganismen drosseln bei extremer Trockenheit trotz hoher Temperaturen ihre Aktivität
In trockenen Perioden (wie im August 2001 beobachtet) bricht die normale
Temperatur-Respirations-Kurve ein. Das Modell für Bodenatmung wird signifikant besser,
wenn sowohl Temperatur als auch Feuchte einbezogen werden

\texttt{temp.15} und \texttt{rootdw0}:
Die Temperatur-Sensitivität (Q10) unterscheidet sich zwischen der autotrophen Respiration (Wurzeln) und der heterotrophen Respiration (Mikroorganismen)
Wurzelrespiration macht oft mehr als die Hälfte der gesamten Bodenatmung aus. Da Wurzeln und
Mikroorganismen unterschiedliche physiologische Optima haben, hängt die Gesamtreaktion des Bodens auf
Erwärmung davon ab, wie viel Wurzelbiomasse vorhanden ist

\texttt{soilc0} und \texttt{rootdw0}
Räumliche Muster der Bodenatmung werden oft gemeinsam durch die Verteilung von Pflanzen (Wurzeln) und Bodenressourcen (C) erklärt
Das beste Modell in der Publikation kombiniert explizit Wurzelbiomasse, Boden-C und Feuchte.
Diese Faktoren korrelieren oft räumlich miteinander, da Bäume über den Wurzeleintrag den Bodenkohlenstoff beeinflussen

\texttt{litter.d} und \texttt{soilc0}:
Die Verfügbarkeit von Substrat (Kohlenstoff) für den Abbau wird durch die Feuchtigkeit limitiert oder gefördert.
Ein sehr hoher Wassergehalt (besonders in tonreichen Böden wie im Hainich) kann die
Sauerstoffdiffusion hemmen und so die mikrobielle Aktivität trotz vorhandenen Kohlenstoffs
einschränken. Zudem sind Streufeuchte (lmoi) und Streudicke (litter.d) zentrale Parameter für die
räumliche Variabilität

\subsection{statistisches Modell}

Zur Vorhersage der Bodenatmung wurde ein lineares Modell (\texttt{lm}) verwendet.\\
Deterministischer Teil: $\mathbb{E}Y_i= \mu_i =\beta_1 + \beta_2x_{i2} + ... + \beta_kx_{ik}$ mit $i= 1,...,n$\\
Stochastischer Teil: $Y_i \sim N(\mu_i, \sigma^2)$ mit $i= 1,...,n$\\
\\
Die globale Suche wurde mit der Funktion \texttt{regsubsets()} durchgeführt. Diese Funktion geht alle möglichen Kombinationen der
Einflussgrößen in einem linearem Modell durch. Für die Qualität der Modelle wurde das BIC genutzt. Der BIC ist das Bayesian Information Criterion und ist
ein Kriterium zur Modellwahl. Es wägt die Anpassungsgüte eines statistischen Modells gegen seine Komplexität, also die Anzahl der
Variablen ab. Das BIC wird berechnet mit der Formel $-2*logL+log(n)*p$ mit der Log-Likelihood (L) des Modells, der Anzahl der
geschätzten Parameter (p) und dem Stichprobenumfang (n). Damit lässt sich einfach der Varianz-Trade off ablesen.

\subsection{Simulationsstudie}

\subsection{Code-Verfügbarkeit}
link zum repository
% ---------------------------------------------------------
\newpage
\section{Ergebnisse}
\textcolor{red}{• Darstellung der Ergebnisse:
• Tabellen
• Grafiken
• Text
• für Datenanalyse und Simulation}

hier tabellen und graphs
\subsection{Modellwahl}

Als Einflussgrößen für das Modell kamen alle Variablen infrage, die eine gewisse Korrelation aufweisen. Als Grenze wurde
eine Korrelation von mehr als 0,2 festgelegt. Damit halbiert sich die Anzahl der Einflussgrößen. Auf den übrigen Variablen
wurde mittels globaler Suche das beste Modell ausgewählt. Ausgewählt wurde dann das Modell mit dem kleinsten BIC.

Bei den Interaktionen wurden alle biologisch plausible ausprobiert. Die Interaktion die den AIC und BIC verbessert hat wurde
benutzt. Durch die Menge der Daten, ist sie nicht statistisch signifikant.

das Modell: $ln(soil.res)=\beta_0+\beta_1*{lmoi}^2+\beta_2*soilc0+\beta_3*({lmoi}^2*soilc0)+\beta_4*{soilc0}^2+\beta_5*{temp.15}^2+\beta_6*{rootdw0}^2+\beta_7*{litter.d}^2+\epsilon$

\subsection{Simulation}

Die Simulationsstudie zeigt, dass ...
% ---------------------------------------------------------
\newpage
\section{Diskussion}
\textcolor{red}{• von der Statistik zur Fragestellung
• Interpretation der statistischen Ergebnisse im Kontext der ursprünglichen Fragestellung
sind die Ergebnisse plausibel/überraschend?
was bedeuten die Ergebnisse
• statistische Diskussion der Simulationsergebnisse}

Die Ergebnisse der Modellwahl zeigen, dass...



% ---------------------------------------------------------
\newpage
\section{Literatur}\label{sec:literatur}
\printbibliography

\end{document}